\documentclass[10pt,CJKmath]{ctexart}
\usepackage{MyNote}



\usepackage{tikz}

% 定义重点高亮和公式框样式
\usepackage{framed,longtable,booktabs,array,calc}

\RequirePackage{tikz-3dplot}






\begin{document}

	\pagestyle{fancynotes}
	\part{Attribute Set}

\section{\texttt{PreAttributeChange}}

\textbf{核心结论：}
$\begin{cases}
	\text{CurrentValue}\Rightarrow\text{PreAttributeChange} \\
	\text{BaseValue}\Rightarrow\text{PreAttributeBaseChange}
\end{cases}$


\begin{enumerate}
	\item \textbf{两者触发时机}

\begin{itemize}
	\item \ci{PreAttributeChange}在属性的 \textbf{当前值（CurrentValue）} 即将变化时触发。  
	常见于带持续时间/无限时长 GE 的修饰器引起的数值变化（聚合后更新 Current）。
	
	\item \ci{PreAttributeBaseChange}
	在属性的 \textbf{基础值（BaseValue）} 即将变化时触发。  
	常见于 Instant GE 或直接改 Base 的路径（如 \texttt{SetNumericAttributeBase}）。
\end{itemize}

\item \textbf{为什么会有两个钩子？}

在 GAS 里通常可理解为：\colorbox{yellow!5}{$
\text{CurrentValue} = \mathbf{f}(BaseValue, Modifiers)
$}
因此：
\begin{itemize}
	\item 改 \textbf{BaseValue} 前，用 \texttt{PreAttributeBaseChange}；
	\item 改 \textbf{CurrentValue} 前，用 \texttt{PreAttributeChange}。
\end{itemize}

\item \textbf{实战建议（很重要）}

\begin{itemize}
	\item 如果你要做夹逼（clamp），例如生命值限制在 \([0, MaxHealth]\)，
	最稳妥做法是 \textbf{两个函数都处理}，避免不同修改路径漏掉。
	\item 对 Instant GE 的“最终结果”校正，通常还会配合下面这个收尾钩子：
	\begin{itemize}
		\item \ci{PostGameplayEffectExecute}
	\end{itemize}
\end{itemize}
\end{enumerate}
在 GAS 中，带 Period 的 GE 每次周期触发会按“执行(Execute)”语义处理，这条路径本质上接近 Instant GE，通常会动到 BaseValue。
而\textbf{一旦 BaseValue 变了，聚合器会重新计算 CurrentValue}，于是你通常会看到：
\begin{itemize}
	\item \ci{PreAttributeBaseChange(...)}：Base 将要改；
	\item \ci{PreAttributeChange(...)}：Current 随后重算将要改。
\end{itemize}
所以“两个都调用”是正常现象，不是异常。

可以把流程理解成：\colorbox{yellow}{$
\text{Periodic Tick} \Rightarrow \Delta \text{BaseValue} \Rightarrow \text{Recompute CurrentValue}$}。

因此：\ci{PreAttributeBaseChange} 先触发，\ci{PreAttributeChange} 后触发。

\paragraph{为什么你会觉得“按道理只改 Current”？}
因为“持续时间/无限时长 GE”这类描述通常让人联想到“临时 Buff/Debuff（Modifier 挂在 Active GE 上）”，那条路径主要表现为 Current 层变化。
但一旦你加了 Period，每个周期会走一次执行逻辑，常常落到 Base 改动路径上，所以就把两边都带起来了。
\begin{lemma}
实战建议（Aura 项目里很实用）
\begin{enumerate}
	\item \ci{PreAttributeBaseChange}：放“基础值合法化”逻辑（如防 NaN、下限保护）。
	\item \ci{PreAttributeChange}：放“显示值 / 运行时值的 clamp”（如 Health 限制到 [0, MaxHealth]）。
	\item \ci{PostGameplayEffectExecute}：处理“执行后副作用”（死亡判定、受击、溢出转护盾等）。
	\item 避免在两个 \ci{Pre} 里写同一套重逻辑，否则容易重复处理。
\end{enumerate}
\end{lemma}
\end{document}